\documentclass{beamer}
\usetheme{Madrid}
\usecolortheme{default}

% Pacchetti necessari
\usepackage[utf8]{inputenc}
\usepackage[italian]{babel}
\usepackage{graphicx}
\usepackage{booktabs}

% Informazioni Titolo
\title[Sentiment Analysis Bi-LSTM]{Sentiment Analysis finanziario tramite Bi-LSTM: pre-train e fine-tuning}
\author[Marchesi, Minei]{Marchesi Luca \and Minei Francesco}
\institute[UniPR]{Università di Parma \\ Laboratorio di Fondamenti di AI}
\date{03 Giugno 2026}

\begin{document}

% Slide del Titolo
\begin{frame}
    \titlepage
\end{frame}

% Slide 1: Obiettivo
\begin{frame}{Obiettivo del progetto}
    \begin{itemize}
        \item \textbf{Task:} Classificazione del sentiment in testi finanziari (Positive, Neutral, Negative).
        \item \textbf{Dataset target:} \textit{Financial PhraseBank} (circa 4.800 frasi etichettate con 50\% agree).
        \item \textbf{Sfida principale:} Scarsità di dati nel dominio specifico e la complessità del linguaggio economico.
        \item \textbf{Strategia:} \textit{Transfer learning e fine-tuning}.
              \begin{itemize}
                  \item Pre-addestramento su \textbf{IMDB Reviews} (50k campioni).
                  \item Fine-tuning sul dominio finanziario.
              \end{itemize}
    \end{itemize}
\end{frame}

% Slide 2: Architettura Base
\begin{frame}{Architettura del modello}
    \begin{block}{Encoder: Bi-LSTM}
        Rete neurale ricorrente bidirezionale per catturare il contesto da entrambi i lati.
    \end{block}
    \begin{block}{Classifier: MLP}
        Testa di classificazione con ReLU e Dropout per la decisione finale.
    \end{block}
    \begin{itemize}
        \item \textbf{Vocabolario:} 20.000 token condivisi tra i due dataset.
        \item \textbf{Embedding:} 100 dimensioni, pre-addestrato su IMDB.
    \end{itemize}
\end{frame}

% Slide 3: Il Problema del Padding Drift
\begin{frame}{Problema tecnico: eccessivo padding}
    \begin{itemize}
        \item Le frasi finanziarie sono brevi ($\sim$20 parole), ma il buffer è fisso a 256 token.
        \item \textbf{Problematica:} L'estrazione dell'ultimo stato nascosto ($h_{256}$) è \textit{contaminata} da centinaia di token di padding.
        \item \textbf{Soluzione:} \textit{Masked Mean Pooling}.
    \end{itemize}
    \begin{exampleblock}{Formula}
        \[ Features = \frac{\sum (h_t \cdot m_t)}{\sum m_t} \]
        Dove $m_t$ è una maschera binaria che ignora i token di padding.
    \end{exampleblock}
\end{frame}

% Slide 4: Strategia di Fine-Tuning
\begin{frame}{Fine-tuning e regolarizzazione}
    Per evitare l'overfitting:
    \begin{enumerate}
        \item \textbf{Layer freezing:} Embedding congelato per preservare le basi linguistiche di IMDB.
        \item \textbf{Learning rate differenti:}
              \begin{itemize}
                  \item LSTM: $10^{-4}$ (adattamento lento).
                  \item MLP: $10^{-3}$ (apprendimento veloce).
              \end{itemize}
        \item \textbf{Ottimizzatore AdamW:} Disaccoppiamento del Weight Decay per regolarizzazione L2 più stabile.
        \item \textbf{Class weights:} Bilanciamento delle classi tramite \textit{Square Root Balancing}.
    \end{enumerate}
\end{frame}

% Slide 5: Analisi Sperimentale - Overfitting
\begin{frame}{Analisi sperimentale e capacità dei modelli}
    \begin{columns}
        \column{0.5\textwidth}
        \textbf{Rete Grande (256x2)}
        \begin{itemize}
            \item 2.44M Parametri.
            \item Train Acc: 99\%.
            \item Val Acc: 69\%.
            \item \alert{Overfitting Critico:} La rete memorizza il training set invece di imparare pattern generali.
        \end{itemize}

        \column{0.5\textwidth}
        \textbf{Rete Semplificata (128x1)}
        \begin{itemize}
            \item 268k Parametri ($-90\%$).
            \item \textbf{Model Squeezing:} La minor capacità forza la rete a generalizzare.
            \item Risultati più stabili e Loss ridotta.
        \end{itemize}
    \end{columns}
\end{frame}

% Slide 6: Risultati Finali
\begin{frame}{Confronto Risultati}
    \begin{table}[]
        \centering
        \begin{tabular}{|l|c|c|c|}
            \hline
            \textbf{Architettura}    & \textbf{Parametri} & \textbf{Val. Acc} & \textbf{Val. Loss} \\
            \hline
            Bi-LSTM (256x2)          & 2.442.243          & 69,51\%           & 0,784              \\
            \textbf{Bi-LSTM (128x1)} & \textbf{268.803}   & \textbf{69.71\%}  & \textbf{0,777}     \\
            \hline
        \end{tabular}
        \label{tab:results}
        \caption{Risultati delle due architetture testate}
    \end{table}
    \begin{itemize}
        \item La versione semplificata ottiene la \textbf{miglior generalizzazione}.
        \item La loss (0.775) dell'architettura semplificata risulta inferiore.
    \end{itemize}
\end{frame}

% Slide 7: Conclusioni
\begin{frame}{Conclusioni e Sviluppi Futuri}
    \begin{itemize}
        \item Risulta evidente che su dataset piccoli architetture meno complesse generalizzano meglio.
        \item Il \textit{Masked Pooling} è essenziale per gestire sequenze di lunghezza variabile.
        \item \textbf{Limite architetturale:} La Bi-LSTM raggiunge un tetto fisiologico del $\sim$70\% su questo dataset di piccole dimensioni.
        \item \textbf{Sviluppi:} Passaggio ad architetture \textit{Transformer} (FinBERT) per catturare relazioni semantiche più complesse.
    \end{itemize}
\end{frame}

\end{document}
